\chapter{Materiais}\label{ch:materials}

Este capítulo descreve a área de estudo, os sete segmentos de uso da terra, a plataforma de computação
e o inventário de dados. As tabelas completas de fontes de dados estão no Anexo~\ref{anx:data}.

\section{Área de estudo}

A área de interesse é o estado de Goiás em conjunto com o Distrito Federal (GO~e~DF)
na região Centro-Oeste do Brasil, cerca de 340.000~km\textsuperscript{2} situados
inteiramente no bioma Cerrado. O Cerrado é uma savana tropical com clima fortemente sazonal
(verão úmido e inverno seco), além de solos profundos, intemperizados, naturalmente ácidos e
de baixa fertilidade (predominantemente Latossolos, ou \emph{Ferralsols}/\emph{Oxisols}) que, no
entanto, podem se tornar altamente produtivos uma vez feita a calagem e a correção de micronutrientes
\citep{Yamada2005}. Nas últimas décadas, o bioma tornou-se um dos principais territórios de grãos,
fibra e pecuária do Brasil --- precisamente o que torna a questão do potencial \emph{versus} uso
atual (QP2) consequente aqui.

A geometria da área de interesse é derivada da Malha Municipal Digital oficial do IBGE (edição 2025,
SIRGAS~2000/EPSG:4674, 246 municípios de Goiás mais o Distrito Federal), dissolvida em um polígono único
que se estende aproximadamente da longitude $-53{,}25^\circ$ a $-45{,}90^\circ$ e da latitude
$-19{,}50^\circ$ a $-12{,}39^\circ$. Essa malha substitui um produto global de limites administrativos
(GAUL) usado numa iteração anterior: o artefato mesclado (\texttt{data/mart/malha\_municipal.gpkg}) é
carregado em memória e convertido em \texttt{ee.FeatureCollection}/\texttt{ee.Geometry} no lado do
cliente a cada execução --- nunca exportado como \emph{asset} persistente do Earth Engine, preservando a
filosofia de zero \emph{upload} (\S\ref{sec:cost}) --- com as geometrias simplificadas em memória
(tolerância $\approx$111~m, cerca de 5$\times$ mais fina que a simplificação de 500~m do GAUL) antes de
compor qualquer requisição ao Earth Engine. Para a fase de uso atual da terra, as próprias 247 unidades
municipais dessa mesma malha --- não mais um produto administrativo global --- alimentam a agregação e
o relato municipal.

\section{Os sete segmentos de uso da terra}

O atlas avalia sete segmentos, escolhidos de forma a cobrir alimento, fibra, proteína animal, energia e
conservação, de modo que a análise possa expressar genuínas compensações (\emph{trade-offs}) entre segmentos.

\begin{table}[H]
\centering
\caption{Os sete segmentos de uso da terra.}\label{tab:segments}
\begin{tabular}{@{}clp{9.5cm}@{}}
\toprule
\# & Segmento & Definição / escopo \\
\midrule
1 & \textbf{Soja} & Soja mecanizada de sequeiro --- a principal cultura comercial da região,
  cuja adaptação agroclimática ao Cerrado está bem caracterizada \citep{Corbellini2024} e que, por si só,
  já se altera sob as projeções de mudança climática \citep{deSouza2026}. \\
2 & \textbf{Cana-de-açúcar} & Cana-de-açúcar mecanizada para açúcar e etanol, com janela de maturação
  seca, seguindo os limiares de temperatura e precipitação estabelecidos para o zoneamento climático da
  cultura no Brasil \citep{Aparecido2021}. \\
3 & \textbf{Outras culturas anuais} & Milho, algodão e sorgo, incluindo a safrinha (segunda-safra). \\
4 & \textbf{Piscicultura} & Aquicultura em tanques (açudes) e reservatórios (espécies de água quente,
  como a tilápia), seguindo critérios de viabilidade local (temperatura da água, textura do solo,
  declive, acesso à água) estabelecidos para aquicultura em viveiros em outras regiões
  \citep{Ssegane2012,Assefa2018}. \\
5 & \textbf{Pecuária} & Gado de corte em pastagem cultivada, que cobre cerca de 28\% do bioma Cerrado
  e abrange um amplo espectro de estados de intensificação e degradação \citep{Vieira2021,Molossi2020}. \\
6 & \textbf{Conservação} & Conservação de floresta e Cerrado --- valor de serviços da vegetação nativa,
  da água e da biodiversidade, ecoando as compensações entre prioridade de restauração e aptidão agrícola
  mapeadas para o bioma \citep{Schuler2022}. \\
7 & \textbf{Geração fotovoltaica} & Geração fotovoltaica em escala de concessionária, com escolha de
  sítio por irradiância, declive, cobertura da terra e acessibilidade, seguindo critérios de seleção de
  local estabelecidos na literatura para usinas fotovoltaicas \citep{Rane2024}. \\
\bottomrule
\end{tabular}
\end{table}

\section{Plataforma de computação e filosofia de custo}\label{sec:cost}

A análise é executada inteiramente no lado do servidor em uma plataforma de observação da Terra na nuvem
(Google Earth Engine), em seu nível gratuito não comercial para pesquisa \citep{Gorelick2017}.
O objetivo desta abordagem é a reprodutibilidade com custo marginal baixo: sem \emph{hardware}
dedicado e sem armazenamento local de imagens de satélite. A máquina local é utilizada apenas para
desenho de gráficos, junção de tabelas e operações simples. Os produtos intermediários são
armazenados na plataforma como imagens raster persistentes, para que etapas subsequentes não
recalculem os grafos anteriores. O processamento permanece dentro da cota de uso, limitando as
amostras de agrupamento e classificação à ordem de $10^5$ pontos.

\begin{table}[H]
\centering
\caption{Matriz de análise e períodos de referência.}\label{tab:grid-periods}
\begin{tabular}{@{}p{3.4cm}p{3.6cm}p{6cm}@{}}
\toprule
Convenção & Valor & Justificativa \\
\midrule
Resolução de análise & 250~m & Corresponde à grade de vegetação do MODIS; situa-se entre a resolução fina de solo/terreno e a grossa do clima. \\
Sistema de referência de coordenadas & EPSG:4326 (geográfico) & A área das células vem da função planimétrica de área de pixel, nunca de um campo armazenado. \\
Normal climatológica & 1991--2020 (padrão da OMM) & Referência única para os deltas de mudança climática. \\
Registro da vegetação & 2001--2020 & Fenologia e indicador de produtividade. \\
Janelas futuras & 2031--2050; 2051--2070 & Duas janelas de projeção. \\
Cenários futuros & SSP2-4.5; SSP5-8.5 & Cenários de emissão moderada e alta. \\
\bottomrule
\end{tabular}
\end{table}

\section{Fontes de dados}

Todos os conjuntos de dados vêm do catálogo público da própria plataforma, acessados na resolução nativa e
convertidos para unidades físicas quando as bandas vêm empacotadas como inteiros. As variáveis
\textbf{biofísicas e agroclimáticas} (Tabela~\ref{tab:src-biophys} do Anexo) cobrem clima e regime de água
(TerraClimate, \citealp{Abatzoglou2018}, com CHIRPS, \citealp{Funk2026}, e WorldClim como suplementos),
terreno (SRTM) e hidrografia (HydroSHEDS). Cobrem também solos (OpenLandMap, 0--30~cm, \citealp{Hengl2026}),
água superficial (JRC Global Surface Water, \citealp{Pekel2016}), fenologia (MODIS NDVI) e acessibilidade
ao mercado (Global Accessibility to Cities). Completam o conjunto a cobertura genérica da terra
(MapBiomas, coleção~10 --- substitui o ESA WorldCover global de 10~m usado numa iteração anterior;
\S\ref{sec:phaseb} justifica o porquê), as áreas protegidas (WDPA), o carbono de biomassa acima do
solo (Spawn \emph{et al.}) e o clima futuro (NASA NEX-GDDP-CMIP6, \citealp{Thrasher2022}).
As entradas de \textbf{uso atual e validação} (Tabela~\ref{tab:src-realized} do Anexo) são as 247
unidades municipais da Malha Municipal Digital do IBGE (edição 2025, substituindo o GAUL usado numa
iteração anterior), as classificações anuais de uso da terra do MapBiomas \citep{Souza2020} e os produtos de
produtividade NPP/GPP do MODIS MOD17 \citep{Zhao2005}. O uso atual da terra e o indicador de produtividade
servem estritamente para mascaramento, validação e caracterização descritiva; nunca entram no cálculo da
viabilidade ou do agrupamento.

\section{Cenários futuros do clima (CMIP6)}

A projeção usa o conjunto NEX-GDDP-CMIP6 da NASA, de resolução temporal diária, em particular os campos de
precipitação e de temperatura máxima e mínima do ar próximo à superfície. Um conjunto de cinco modelos ---
ACCESS-CM2, EC-Earth3, GFDL-ESM4, MPI-ESM1-2-HR e MRI-ESM2-0 --- é usado como subconjunto padrão e
bem-comportado de modelos de circulação geral (GCMs); todos os cinco têm dados para os dois cenários de
interesse. A projeção cobre os cenários SSP2-4.5 e SSP5-8.5, para as janelas de 2031--2050 e 2051--2070.

Relativamente à base histórica de cada modelo, que cobre 1991--2014, a precipitação entra como fator de
mudança multiplicativo e a temperatura, como fator aditivo (em kelvin). A especificação completa está na
Tabela~\ref{tab:cmip6-spec} do Anexo.


\section{Ambiente computacional}

A implementação é orientada a configuração: a lógica analítica é modular e todo parâmetro numérico ---
limiares de pertinência fuzzy, pesos dos fatores, identificadores de conjuntos de dados e fatores de
conversão de unidades --- é declarado como valor versionado, em vez de embutido no código, de modo que a
análise possa ser auditada e reexecutada apenas a partir da especificação declarada. O software é
deliberadamente simples: além da integração com o Google Earth Engine e da orquestração das etapas de
análise, o código Python usa bibliotecas científicas comuns para as pequenas tarefas executadas
localmente. Nenhum arcabouço de aprendizado profundo é utilizado.

O código-fonte completo --- configurações, módulos de análise, cadernos e ferramentas de exportação
--- está disponibilizado publicamente no repositório \texttt{goias\_suitability}
\citep{daSilva2026}. O atlas interativo construído a partir destes resultados está publicado em
\url{https://ahxsahxs.github.io/goias_suitability}.